\documentclass{subfiles}
\begin{document}
    \begin{definition*}    
        Let \(X\) be a topological space. We say that \(X\) is Noetherian 
            if and only if, any descending chain of closed sets 
            \[
                X = X_{1} \supseteq X_{2} \cdots \supseteq X_{k} \supseteq  \cdots
            \]
            is stationary.
    \end{definition*}

    Note that if \(X \subseteq \mathbb{A}^{n}\), then \(X\) is a Noetherian topological space.
        We now recall what a reducible/inreducible space is.
    \begin{definition*}
        Let \(X\) be a topological space. Then, 
            \(X\) is reducible if \(X = X_{1} \cup X_{2}\) with \(X_{1}, X_{2}\)
            proper closed subsets. Is inreducible viceversa.
    \end{definition*}

    \begin{lemma*}
        Let \(X = X_{1} \cup \cdots \cup X_{k}\) be a topological space,
            where \(X_{i}\) is a proper closed set. Then, \(X\) is reducible.
    \end{lemma*}

    \begin{proof*}
        Let us prove the result by induction on the number of closed sets.
        \begin{description}
            \item [\(k = 2\):] follows by the definition.
            \item [\(k \ge 3\):] Assume that for \(k - 1\) it holds that \(X\)
                is reducible. This means that \(X = X_{1} \cup \cdots \cup X_{k - 1}\).
                But each \(X_{i}\) is a closed set; additionally, a finite union of 
                closed sets is still a closed set. Hence, we have that 
                \[
                    X = \underbrace{Z}_{X_{1} \cup \cdots \cup X_{k - 1}} \cup X_{k}.
                \]
                Therefore \(X\) is reducible.\qed
        \end{description} 
    \end{proof*}

    \begin{lemma*}
        Let \(X\) be a topological space. Then, the following are equivalent.
        \begin{enumerate}
            \item \label{case:1} \(X\) is inreducible.
            \item \label{case:2} For any \(U_{1}, U_{2} \subseteq X, U_{i} \ne \varnothing\), 
                then \(U_{1} \cap U_{2} \ne \varnothing\).
            \item \label{case:3} Each non-empty open set \(U \subseteq X\) is dense in \(X\).
        \end{enumerate}
    \end{lemma*}

    \begin{proof*}
        \begin{description}
            \item [\(\ref{case:1} \implies \ref{case:2}\):] Let \(U_{1}, U_{2}\) be non-empty open sets 
                of \(X\). We can write 
                \[
                    U_{1} = X \setminus X_{1} \text{ and } \quad U_{2} = X \setminus X_{2},
                \]
                with \(X_{1}, X_{2}\) proper closed sets of \(X\).
                By the sake of contradiction, let \(U_{1} \cap U_{2} = \varnothing\).
                Then we have that 
                \[\begin{aligned}
                    U_{1} \cap U_{2} & = X \setminus X_{1} \cap X \setminus X_{2} \\
                        & = X \setminus (X_{1} \cup X_{2}) = \varnothing.
                \end{aligned}\]
                Thus we have \(X = X_{1} \cup X_{2}\), which is a contradiction.

            \item [\(\ref{case:2} \implies \ref{case:3}\):] Let \(U \subseteq X\) be an non-empty open set.
                Then \(\forall W \subseteq X\), such that \(W\) is open and 
                non-empty, we have that \(U \cap W \ne \varnothing\). 
                By the criterion for density, it follows that \(U\) is dense.

            
            \item [\(\ref{case:3} \implies \ref{case:1}\):] Let check this one. % Assume \(X = X_{1} \cup X_{2}\) with \(X_{i}\)
            %     proper closed sets. Let \(U_{1} = X \setminus X_{1}, U_{2} = X \setminus X_{2}\)
            %     be non-empty open sets. By 2., we should 
                \qed
        \end{description}
    \end{proof*}
    \clearpage

    \begin{theorem*}
        Let \(V \subseteq \mathbb{A}^{n}\) be an affine closed set. Then,
        \(V\) is inreducible if and only if \(I(V)\) is a prime ideal.
    \end{theorem*}

    \begin{proof*}
        \begin{description}
            \item [\(\impliedby \):] By the sake of contradiction let \(V\) be reducible. 
                Then, \(V = V_{1} \cup V_{2}\), with \(V_{i}\) be proper closed sets.
                We have that 
                \[
                    V_{1} \subset V \text{ and } V_{2} \subset X.
                \]
                Considering the ideals, we have 
                \[
                    I(V_{1}) \supset I(V) \text{ and } I(V_{2}) \supset I(V).
                \]
                Hence, there exist \(F_{1}, F_{2} \in \K\) such that \(F_{1} \in I(V_{1}),
                F_{2} \in I(V_{2}), F_{1}, F_{2} \ne I(V)\). Thus, 
                \(\forall P \in V, F_{1}F_{2}(P) = F_{1}(P)F_{2}(P) = 0\).
                Meaning \(P \in V_{1} \lor P \in V_{2} \implies F_{1}F_{2} \in I(V)\)
                which is a contradiction.

            \item [\(\implies\):] Assume that \(I(V)\) is not prime.
                Therefore, \(\exists F_{1}, F_{2} \in \K, F_{1},F_{2} \notin I(V) \land F_{1}F_{2} \in I(V)\).
                
                Let \(S_{i} = I(V) \cup \set{F_{i}}\). We have 
                \[\begin{aligned}
                    V(S_{i}) &= V(\set{F_{i}} \cup I(V)) \\ 
                     &= V(\set{F_{i}}) \cap V, i = 1, 2 \\ 
                     &= \set{P \in V}[F_{i}(P) = 0].
                \end{aligned}\]
                Thus, each \(V_{i}\) is a proper closed subset of \(V\).

            We prove that \(V_{1} \cup V_{2} = V\). Let \(P \in V\).
                Then, \(F_{1}F_{2}(P) = 0 \implies P \in V_{1} \lor P \in V_{2}\),
                hence \(P \in V_{1} \cup V_{2}\).\qed
        \end{description}
    \end{proof*}

    \begin{remark*}
        Note that if \(\F\) is an infinite field, then \(\mathbb{A}^{n}\) is inreducible.

        In fact, if \(V = \mathbb{A}^{n}\), Considering the ideals this generate we have 
        \[
            I(V) = I(\mathbb{A}^{n}) = \set{F \in \K[x_{1}, \ldots, x_{n}]}[F(P) = 0, \forall P \in \mathbb{A}^{n}].
        \]
        By the identity principle \todo{Search what this ``indentity principle'' is.}, we have that the only
        polynomial in \(I(V)\) is the null one; hence \(I(V)\)  is an integral domain,
        thus is prime and therefore \(V\) is inreducible.
    \end{remark*}

    \begin{proposition*}
        Noetherian topological spaces satisfy the following.
        \begin{enumerate}
            \item If \(X\) is a Noetherian topological space and \(Y \subseteq X\) 
                is a topological subspace, then \(Y\) is Noetherian.

            \item The union of finitely many Noetherian topological spaces is Noetherian.
            \item If \(X\) is a Noetherian topological space, then \(X\) is semi-compact.
        \end{enumerate}
    \end{proposition*}

    \begin{proof*}
        \begin{enumerate}
            \item Note that we can write \(Y\) as a chain of closed sets; that is,
                \(Y = Y_{1} \supseteq Y_{2} \cdots \supseteq Y_{k} \supseteq \cdots\),
                where each \(Y_{i} = Y \cap X_{i}\). What we are left with is proving that the 
                \(X_{i}\)s form a chain in \(X\). Let us rewrite the \(X_{i}\).

                Let \(X_{1}' = X_{1}\). For any \(i \ge 2\) we have \(Y_{i} = Y \cap X_{i}\)
                and \(Y_{i} \subseteq Y_{i - 1}\), meaning that \(Y_{i} \subseteq X_{i - 1}\);
                that is, \(Y_{i} = Y \cap (X_{i} \cap \cdots \cap X_{1}) = Y \cap X_{i}'\).

                From the above \(X_{1}' \subseteq X_{2}' \subseteq \cdots \subseteq X_{k}' \subseteq \cdots\).
                Since \(X\) is Noetherian, the chain of \(X_{i}'\) must be stationary,
                hence the chain of \(Y_{i}\) is also stationary.

            \item Since we are considering a finite union, we can simply prove the 
                proposition for \(n = 2\), as the proof for \(n \ge 3\) follows by induction.

                Let \(Y, Z\) be Noetherian spaces. We prove that \(Y \cup Z = X\) is a Noetherian
                    spaces. 

                Since \(Y, Z\) are Noetherian we know that 
                \[\begin{aligned}
                    Y & = Y_{1} \supseteq Y_{2} \supseteq \cdots \supseteq Y_{k} \supseteq \cdots \\ 
                    Z & = Z_{1} \supseteq Z_{2} \supseteq \cdots \supseteq Z_{l} \supseteq \cdots. \\ 
                \end{aligned}\]
                Additionally, each \(Y_{i} \text{ and } Z_{i}\) can be seen as the intersection of some 
                    closed set of \(X\) with \(Y \text{ and } Z\) respectively.
                    That is, \(Y_{i} = Y \cap X_{i}\) and \(Z_{i} = Z \cap X_{i}\).

                Now we know that \(\exists s \in \N\) such that \(Y_{s} = Y_{s + 1} = \cdots\) and 
                    \(\exists t \in \N\) such that \(Z_{t} = Z_{t + 1} = \cdots\). Hence, 
                    if \(m = \max{s, t}\) we have that \(Y_{m} = Y_{Y + 1} = \cdots\) and 
                    \(Z_{m} = Z_{m  + 1} = \cdots\). But for all \(i\), it holds 
                    \[
                        X_{i} = (Y \cap X_{i}) \cup (Z \cap X_{i}) = (Y \cup Z) \cap X_{i}
                    \]
                    by which \(X_{m} = X_{m + 1} = \cdots\), meaning that \(X\) is Noetherian.

                \item Let \(X = \bigcup_{i \in I} U_{i}\), with \(I\) set of indexes.

                    Let \(i_{1} \in I\). If \(X = U_{i_{1}}\) we are done.
                    Otherwise, there exists some \(i_{2}\) such that \(U_{i_{2}} \subset U_{i_{2}}\).
                    If \(U_{i_{1}} \cup U_{i_{2}} = X\) we are done, otherwise we repeat the process.

                    Two scenarios follow:
                    \begin{enumerate}
                        \item \(\exists i_{n} \in I\) such that \(U_{i_{1}} \cup U_{i_{2}} \cup \cdots \cup U_{i_{n}} = X\).
                            In which case we get the thesis; or 
                        \item We get a chain of open sets that is not stationary. Meaning that,
                            once we consider the complements, we have a chain of closed sets that is
                            not stationary, which contradicts the hypotesis of \(X\) being Noetherian.\qed
                    \end{enumerate}
            \end{enumerate}
    \end{proof*}


    \begin{lemma*}
        Let \(X\) be a Noetherian space. If \(\F* \ne \varnothing\) is a family of closed sets,
        then exists \(m \in \F*\) minimum.
    \end{lemma*}

    \begin{proof*}
        Trivial. Analogous to the proof of the existence of a maximal element in a Noetherian ring.\qed
    \end{proof*}

    \subsubsection{Unique decomposition}
    \subfile{../subsub/2.1.1 - Unique decomposition}

\end{document}
